\section{Замкнутость регулярных языков относительно теоретико-множественных операций}
\label{sec:RegularClosureProperties}
\tikzsetfigurename{RegLang_Closure_}

\begin{theorem}
    Класс регулярных языков над алфавитом $\Sigma$ замкнут относительно следующих теоретико-множественных операций:
    \begin{enumerate}
        \item Объединение: если $L_1, L_2$~--- регулярные, то $L_1 \cup L_2$~--- регулярный;
        \item Пересечение: если $L_1, L_2$~--- регулярные, то $L_1 \cap L_2$~--- регулярный;
        \item Дополнение: если $L$~--- регулярный, то $\overline{L}$~--- регулярный;
        \item Разность: если $L_1, L_2$~--- регулярные, то $L_1 \setminus L_2$~--- регулярный.
    \end{enumerate}
\end{theorem}

\begin{proof}
    Пусть $L_1$ и $L_2$ распознаются конечными автоматами над одним алфавитом $\Sigma$:
    \[
        M_1 = \langle Q^1, Q_S^1, Q_F^1, \delta^1, \Sigma \rangle, \qquad
        M_2 = \langle Q^2, Q_S^2, Q_F^2, \delta^2, \Sigma \rangle.
    \]


    \paragraph{1. Объединение.}
    Построим недетерминированный автомат с $\varepsilon$-переходами,
    вводя новое начальное состояние $q_0$, из которого по $\varepsilon$ можно попасть
    в начальные состояния обоих автоматов:
    \[
        M = \langle Q^1 \cup Q^2 \cup \{q_0\},\; \{q_0\},\; Q_F^1 \cup Q_F^2,\; \delta,\; \Sigma \rangle,
    \]
    где $\delta$ совпадает с $\delta^1$ и $\delta^2$ на соответствующих состояниях, а дополнительно
    \[
        \delta(q_0, \varepsilon) = Q_S^1 \cup Q_S^2.
    \]
    Тогда $L(M) = L_1 \cup L_2$.

    \paragraph{2. Пересечение.}
    Пусть автоматы детерминированы. Построим автомат произведения:
    \[
        M_3 = \langle Q^1 \times Q^2,\; Q_S^1 \times Q_S^2,\; Q_F^1 \times Q_F^2,\; \delta^3,\; \Sigma \rangle,
    \]
    где переходы определяются как:
    \[
        \delta^3\big((q^1, q^2), t\big) = \big(\delta^1(q^1, t),\; \delta^2(q^2, t)\big).
    \]
    Тогда%
    \sidenote{Аккуратно доказать это можно расписав переходы между конфигурациями соответствующих автоматов при обработке цепочки.}
    \[
        w \in L(M_3) \iff w \in L_1 \text{ и } w \in L_2,
    \]
    следовательно, $L(M_3) = L_1 \cap L_2$.

    \paragraph{3. Дополнение.}
    Пусть $M = \langle Q, \{q_S\}, Q_F, \delta, \Sigma \rangle$~--- детерминированный полный автомат.
    Построим автомат
    \[
        M' = \langle Q, \{q_S\}, Q \setminus Q_F, \delta, \Sigma \rangle.
    \]
    Очевидно, $L(M') = \overline{L(M)}$.

    \paragraph{4. Разность.}
    Разность выражается через пересечение и дополнение:
    \[
        L_1 \setminus L_2 = L_1 \cap \overline{L_2}.
    \]
    Регулярные языки замкнуты относительно обеих операций, следовательно, и относительно разности тоже замкнуты.
\end{proof}

%\begin{example}[Конструкция произведения]\label{ex:product-automaton}
%    Пусть $M_1$ и $M_2$~--- детерминированные автоматы над алфавитом $\Sigma=\{a,b\}$.
%    В автомате $M_1$ есть переходы
%    $q^1 \xrightarrow{a} p^1$, $q^1 \xrightarrow{b} p^1$ и $p^1 \xrightarrow{b} p^1$,
%    а в автомате $M_2$~--- переходы
%    $q^2 \xrightarrow{a} q^2$ и $q^2 \xrightarrow{b} p^2$.
%    Тогда в автомате произведения $M_1 \times M_2$ возникают переходы:
%    \begin{align*}
%        (q^1, q^2) \xrightarrow{a} (p^1, q^2), \quad
%        (p^1, q^2) \xrightarrow{b} (p^1, p^2), \quad
%        (q^1, q^2) \xrightarrow{b} (p^1, p^2),
%    \end{align*}
%    что проиллюстрировано на рис.~\ref{fig:product-automaton-example}.
%    Принимающими являются пары из принимающих состояний обоих автоматов.
%\end{example}
%
%
%\begin{marginfigure}
%    \centering
%    \scalebox{0.64}{
%    \begin{tikzpicture}
%        \node[elliptic state,initial] (q0) {$(q^1, q^2)$};
%        \node[elliptic state] (p0) [right=of q0] {$(p^1, q^2)$};
%        \node[elliptic state,accepting] (p1) [below right=of p0] {$(p^1, p^2)$};
%
%        \path[->]
%            (q0) edge[bend left, above] node {$a$} (p0)
%            (p0) edge[bend left, right] node {$b$} (p1)
%            (q0) edge[bend right, below left] node[swap] {$b$} (p1);
%    \end{tikzpicture}}
%    \caption{Пример переходов в произведении автоматов.}
%    \label{fig:product-automaton-example}
%\end{marginfigure}
%


%\section{Вопросы и задачи}
%
%Построить базу.
%
%Научиться выполнять запросы через линейку.
